\subsection{The three framework converses with their exact host}

\begin{hbdefinition}{def:p12-har-fw-host}{The beta-only LF presentation used by the three signatures}
For the framework signatures in
\cref{def:p12-har-prop-signature,def:p12-har-pred-signature,def:p12-har-lambda-signature},
use weak dependent LF with beta conversion, opaque signature constants,
and object-variable contexts. Its split syntax has kinds
\(K::=\mathrm{type}\mid\Pi x:B.K\), families
\(B::=a\mid\Pi x:B.C\mid\lambda x:B.C\mid BM\), and
objects \(M::=c\mid x\mid\lambda x:B.M\mid MN\).
Here \(a\) is a family constant and \(c\) an object constant.
In the PTS view the distinguished kind \(\mathrm{type}\) is
\(\ast\), the top classifier is \(\Box\), the axiom is
\(\ast:\Box\), and the product triples are exactly
\((\ast,\ast,\ast)\), \((\ast,\Box,\Box)\).
The seven underlying PTS rules are the already registered
\hyperref[sheet:p12-har-pts]{Figure~3 surface}, at this specification.
Treat signature declarations as fixed opaque entries when using that view;
do not bind or substitute them as additional object variables.

Thus \(\mathrm{prop}\), \(A\), \(D\) are fixed family constants,
and \(T:\mathrm{prop}\to\mathrm{type}\) and
\(\mathrm{eq}:D\to D\to\mathrm{type}\) are families, not
proposition-valued predicates inside the direct lambdaHOL system.
There is no kind-domain product pair, equality reflection, inductive
eliminator or additional proof-producing axiom. This is the beta-only
LF interface discussed by the primary framework account, not its
separately discussed eta extension and not the three-sort lambdaHOL
\cite[pp.~1160--1163]{BarendregtGeuvers2001}
\cite[sections~2.1 and 2.4]{HarperHonsellPlotkin1993LF}.
\end{hbdefinition}

\begin{hblemma}{thm:p12-har-fw-normalization}{proved=annex:p12.har.fw-normalization}{The exact LF host has typed beta-normal representatives}
Every derivable judgment over a well-formed signature/context in this
host has a beta-normal subject of the same type. The identity sort/rule
inclusion in exact CC gives strong normalization; the source rule
construction gives LF subject reduction. Normal-form generation exposes
the head constant/variable and its complete typed application spine.
Opaque atomic families \(Tq\), \(\mathrm{eq}\,uv\),
\(\mathrm{prop}\), \(A\), and \(D\) cannot become dependent
products by beta conversion. This is not an assumption that every
normal term is beta-convertible to an eta-long canonical term.
\end{hblemma}

\begin{hbdefinition}{def:p12-har-fw-contexts}{Admissible contexts and pointwise reification}
For PROP, the context after \(\Sigma_{\mathrm{PROP}}\) consists
of finitely many distinct atom names \(a_i:\mathrm{prop}\) and
proof variables \(h_j:T(q_j)\). For PRED, after the full signature
(including inherited \(T\)), it consists of object variables
\(x_i:A\) and proof variables \(h_j:T(q_j)\).
The \(q_j\) are well-typed formula names, normalized when decoding.
For the untyped-lambda converse, the context after
\(\Sigma_{\mathrm{lambda}}\) consists only of free object
variables \(x_i:D\); an extra equality-proof assumption would assert
an extra equation, not pure beta equality.

Normalize before reification. PROP formula names are atoms or
\(\mathrm{imp}\,p q\). PRED object names are variables, \(c\),
or \(ft\); formula names are \(Rst\), \(\mathrm{imp}\,p q\),
or \(\mathrm{forall}\,F\). Decode the first two formula cases
structurally. For the third, decode
\(\mathrm{forall}\,F\) as
\(\forall z:A.\mathcal D(\mathrm{nf}_\beta(Fz))\), with \(z\)
fresh. A normal \(F:A\to\mathrm{prop}\) is either
\(\lambda z:A.q\) or \(Rs\); the latter case directly yields
\(\forall z:A.R(\mathcal D(s),z)\).

For normal untyped-lambda names, decode \(x:D\) to the source
variable, \(\mathrm{app}\,uv\) to the source application, and
\(\mathrm{abs}\,F\) to
\(\lambda z.\mathcal D(\mathrm{nf}_\beta(Fz))\).
A normal \(F:D\to D\) is either \(\lambda z:D.u\) or
\(\mathrm{app}\,u\); the latter directly yields
\(\lambda z.\mathcal D(u)z\). These pointwise clauses handle
eta-short LF functions without adding an object-language eta equation.
For example \(\mathrm{abs}(\mathrm{app}\,u)\) decodes to
\(\lambda z.\mathcal D(u)z\), not to \(\mathcal D(u)\).
Freshness is relative to the entire current context and the arguments.
\end{hbdefinition}

\begin{hblemma}{thm:p12-har-fw-reification}{proved=annex:p12.har.fw-reification}{The reifiers are defined and commute with open substitution}
In the stated contexts, the normal-name grammars above are exhaustive,
independent of proof variables, and their reifiers terminate. Reification
of the source encoding recovers each source expression up to bound-name
renaming. Reification is invariant under LF beta conversion and commutes
with substitutions of object names for object variables, including
capture-avoiding substitution under \(\mathrm{forall}\) and
\(\mathrm{abs}\). In PRED these are arbitrary first-order term
substitutions, not assignments only to closed constants; in the lambda
case they are arbitrary finite open untyped-term substitutions.
\end{hblemma}

\begin{hbtheorem}{thm:p12-har-fw-faithfulness}{proved=annex:p12.har.fw-faithfulness}{Derivability reflection for all three source framework signatures}
At exactly the host/signature/context discipline above:
\begin{enumerate}
 \item If \(\Sigma_{\mathrm{PROP}},\vec a,\vec h
 \vdash p:T([\varphi])\), then the decoded assumptions
 \(\mathcal D(q_j)\) derive \(\varphi\) in minimal implication
 natural deduction. With no \(\vec h\), \(\varphi\) is a theorem.
 \item If \(\Sigma_{\mathrm{PRED}},\vec x,\vec h
 \vdash p:T([\varphi])\), then those decoded assumptions derive
 \(\varphi\) in the specified one-sorted minimal
 \(\{\supset,\forall\}\)-logic with \(c,f,R\).
 The constructed universal-introduction rule obeys its eigenvariable
 condition, and elimination uses the actual decoded open argument.
 \item If \(\Sigma_{\mathrm{lambda}},\vec x:D
 \vdash p:\mathrm{eq}[P][Q]\), then \(P=_\beta Q\)
 in the untyped lambda calculus, with its actual open free variables.
\end{enumerate}
The proof constructs a source derivation from every normal proof head,
including function-valued proof arguments to introduction/xi constants.
Combined with \cref{thm:p12-har-framework-forward}, these give both
existence directions. They do not assert a bijection on all LF terms or
faithfulness after adding arbitrary signature axioms
\cite[pp.~1161--1163]{BarendregtGeuvers2001}.
\end{hbtheorem}

\begin{hbcorollary}{thm:p12-har-fw-no-eta}{proved=annex:p12.har.fw-no-eta}{The lambda signature has not acquired object eta}
In the pure lambda signature with \(u:D\), no term inhabits
\(\mathrm{eq}\,u\,[\lambda x.ux]\), for \(x\) fresh.
Pointwise reification of LF function arguments is not the object
equation \(u=\lambda x.ux\).
\end{hbcorollary}

\subsection{The source rational subsetoid construction completed}

\begin{hbtheorem}{thm:p12-har-rational-normalization}{proved=annex:p12.har.rational-normalization}{Reduced representatives and the actual subsetoid isomorphism}
Use the source carrier \(A=\mathbb Z\times\mathbb N\), with
\((a,p)=_A(b,q)\) when \(a(q+1)=b(p+1)\), and
\(P(a,p)\) when \(\gcd(|a|,p+1)=1\). Its subsetoid
\(S=\Sigma z:A.Pz\) compares first components by \(=_A\),
ignoring the proof witnesses. Every class has exactly one coprime
integer-numerator/positive-denominator pair. The construction is
\[
 d=\gcd(|a|,p+1),\qquad
 r(a,p)=\left(a/d,(p+1)/d-1\right),
\]
together with the proved coprimality witness. The first-component
inclusion \(j:S\to A\) and this \(r:A\to S\) respect the
declared equalities, and \(jr\) and \(rj\) are pointwise equal
to their respective identities. This is a setoid isomorphism, not
definitional equality of all proof-witness terms. Zero has reduced pair
\((0,0)\), and negative numerators require no sign change of the
positive denominator \cite[p.~1173]{BarendregtGeuvers2001}.
\end{hbtheorem}

\begin{hbtheorem}{thm:p12-har-setoid-quotient-universal}{proved=annex:p12.har.setoid-quotient-universal}{The full quotient-setoid mapping property}
For a source setoid \([A,=_A]\) and equivalence relation \(Q\)
with \(=_A\subseteq Q\), the identity underlying map
\(q:[A,=_A]\to[A,Q]\) is equality preserving. Maps from
\([A,Q]\) to \([B,=_B]\) correspond precisely to source
setoid maps from \([A,=_A]\) that preserve \(Q\).
The descended map is the same underlying function and is unique up to
the source's pointwise function-setoid equality. No quotient selector,
choice axiom or host proof-irrelevance principle is required.
\end{hbtheorem}

\begin{proofboundary}{def:p12-har-fw-scope-boundary}{The equality and completeness regimes remain separate}
The lambda framework's \(\mathrm{eq}\) is a type representing
object beta equality. It is not host definitional equality, the
source's impredicative Leibniz equality, or a proof-dependent identity
eliminator. The three reflection theorems do not supply the source-cited
global extensional/setoid interpretation for every dependent constructor,
or the separately attributed weak completeness results for Martin--Lof
presentations. Those need their actual type, transport, universe and
elimination hypotheses. The rational and quotient arguments here close
the stated concrete setoid constructions only; no larger type theory
is identified with them.
\end{proofboundary}
